% 谷神星（综述）
% license CCBYSA3
% type Wiki

本文根据 CC-BY-SA 协议转载翻译自维基百科\href{https://en.wikipedia.org/wiki/Ceres_(dwarf_planet)}{相关文章}。

\begin{figure}[ht]
\centering
\includegraphics[width=6cm]{./figures/ed4afc93bd247d5b.png}
\caption{Dawn 探测器于 2015 年 5 月拍摄的 Ceres 图像。两个明亮斑点分别为 Haulani 撞击坑（右）和 Oxo 撞击坑坑底（左）。} \label{fig_Ceres_1}
\end{figure}
Ceres（小行星编号：1 Ceres）是一颗位于 Mars 与 Jupiter 轨道之间主小行星带中的矮行星。它是人类发现的首颗小行星，于 1801 年 1 月 1 日由 Giuseppe Piazzi 在 Sicily 的 Palermo Astronomical Observatory 发现，并在当时被宣布为一颗新行星。此后，Ceres 被重新分类为小行星，近年又被认定为矮行星——它是唯一一颗轨道未超出 Neptune 之外的矮行星，也是所有不具备卫星的矮行星中最大的一颗。

Ceres 的直径约为 Moon 的四分之一。由于其体积很小，即使在最亮时，除非在极其黑暗的天空条件下，也无法用肉眼看到。其视星等范围为 6.7 至 9.3，并在每隔 15 至 16 个月发生一次会合（即靠近 Earth）时达到峰值。因此，即使使用最强大的望远镜，其表面特征也只能勉强识别，在 NASA Dawn 无人探测器于 2015 年抵达并进入 Ceres 轨道任务之前，人们对其所知甚少。

Dawn 的观测发现，Ceres 的表面由水、冰以及碳酸盐和黏土等含水矿物混合组成。重力数据表明，Ceres 的内部可能呈部分分化结构，拥有一个泥状（冰—岩）地幔/核心，以及一层密度较低但更坚固的地壳（其体积最多约 30\% 为冰）。虽然 Ceres 可能缺乏内部液态水海洋，但盐水仍在外地幔中流动并抵达表面，从而大约每五千万年形成一次如 Ahuna Mons 这样的低温火山。这使 Ceres 成为距离 Sun 最近的已知低温火山活跃天体。Ceres 具有极其稀薄且短暂的水蒸气大气，由其表面局部区域喷发形成。
\subsection{历史}  
\subsubsection{发现}  
在 18 世纪日心说被接受到 1846 年 Neptune 被发现之间的多年里，一些天文学家认为数学定律预测在 Mars 和 Jupiter 轨道之间存在一颗隐藏或失踪的行星。1596 年，理论天文学家 Johannes Kepler 认为，只有在添加两颗行星——一颗位于 Jupiter 与 Mars 之间，另一颗位于 Venus 与 Mercury 之间——时，行星轨道之间的比例才会符合“上帝的设计”\(^{[19]}\)。其他理论家，例如 Immanuel Kant，思考该空隙是否由 Jupiter 的引力造成；1761 年，天文学家兼数学家 Johann Heinrich Lambert 问道：“谁又能知道是否已有行星从 Mars 与 Jupiter 之间的广阔空间中消失？天体是否如同 Earth 一样，强者磨损弱者，Jupiter 与 Saturn 是否注定永远掠夺它们？”\(^{[19]}\)

1772 年，德国天文学家 Johann Elert Bode 引用 Johann Daniel Titius 的研究，发表了一条后来被称为 Titius–Bode 定律的公式，该公式似乎能预测已知行星的轨道，但在 Mars 与 Jupiter 之间出现无法解释的空隙\(^{[19][20]}\)。该公式预测 Sun 附近约 2.8 AU（约 4.2 亿 km）处应存在另一颗行星\(^{[20]}\)。随着 William Herschel 于 1781 年在预测距离附近发现 Saturn 之外的 Uranus，这一定律获得更多信任\(^{[19]}\)。1800 年，由德国天文期刊《Monthly Correspondence》主编 Franz Xaver von Zach 领导的一组人向 24 名经验丰富的天文学家发出请求，将他们称为“天上警察”\(^{[20]}\)，希望他们共同努力，开始系统搜索这颗预期行星\(^{[20]}\)。尽管他们未发现 Ceres，但后来发现了小行星 Pallas、Juno 和 Vesta\(^{[20]}\)。

受邀参与该搜索计划的天文学家之一是 Giuseppe Piazzi，他是 Sicily Palermo 学院的一名天主教神父。在收到邀请函之前，Piazzi 于 1801 年 1 月 1 日发现了 Ceres\(^{[21]}\)。他当时在寻找“la Caille 氏黄道星表中的第 87 颗恒星”，却发现“前面还有另一颗”\(^{[19]}\)。Piazzi 并非找到了一颗恒星，而是一颗运动中的类星体，他起初认为那是一颗彗星\(^{[22]}\)。Piazzi 共观测 Ceres 24 次，最后一次观测发生在 1801 年 2 月 11 日，后因生病中断。他于 1801 年 1 月 24 日通过信件向两位天文学家宣布了他的发现：一位是来自 Milan 的同胞 Barnaba Oriani，另一位是 Berlin 的 Bode\(^{[23]}\)。他将其报告为彗星，但表示“由于它的运动如此缓慢且近乎均匀，我多次想到它可能比彗星更重要”\(^{[19]}\)。4 月，Piazzi 将完整观测结果送给 Oriani、Bode 和法国天文学家 Jérôme Lalande。这些信息发表于 1801 年 9 月的《Monatliche Correspondenz》\(^{[22]}\)。

当时，Ceres 的视位置已经发生改变（主要由于 Earth 绕 Sun 运动），且过于靠近 Sun 的眩光，其他天文学家无法确认 Piazzi 的观测。到了年底，Ceres 应该能再次被观测到，但由于已经过去很长时间，其位置难以预测。为重新定位 Ceres，24 岁的数学家 Carl Friedrich Gauss 开发了一套高效的轨道确定方法\(^{[22]}\)。他在几周内预测出 Ceres 的轨迹，并将结果发送给 von Zach。1801 年 12 月 31 日，von Zach 与同为“天上警察”的 Heinrich W. M. Olbers 在预测位置附近再次发现 Ceres，并继续记录其位置\(^{[22]}\)。Ceres 与 Sun 相距 2.8 AU，几乎完美符合 Titius–Bode 定律；但当 1846 年 Neptune 被发现时，其实际位置比预测近 8 AU，多数天文学家认为该定律仅是巧合\(^{[24]}\)。

早期观测者只能大致估算 Ceres 的大小。Herschel 在 1802 年低估其直径为 260 km（160 mi）；1811 年，德国天文学家 Johann Hieronymus Schröter 将其高估为 2,613 km（1,624 mi）\(^{[25]}\)。20 世纪 70 年代，红外光度测量使其反照率测定更为精确，Ceres 的直径得出与真实值 939 km（583 mi）相差不到 10\% 的结果\(^{[25]}\)。
\subsubsection{名称与符号}  
Piazzi 为其发现所提出的名称是 Ceres Ferdinandea：“Ceres”取自罗马农业女神，其神庙最早建立于 Sicily；“Ferdinandea”则是为了纪念 Piazzi 的君主与资助者——Sicily 的 Ferdinand III 国王\(^{[22]}\)。后一个名称未被其他国家接受，因而被取消。在 von Zach 于 1801 年 12 月重新观测到 Ceres 之前，von Zach 将该行星称为 Hera，而 Bode 则称其为 Juno。尽管 Piazzi 表示反对，这些名称在德国一度通行，直到该天体的存在得到确认。确认之后，天文学家最终采用了 Piazzi 所命名的名称\(^{[26]}\)。

Ceres 的形容词形式有 Cererian\(^{[27][28]}\) 和 Cererean\(^{[29]}\)，两者读音均为 /sɪˈrɪəriən/\(^{[30][31]}\)。1803 年发现的一种稀土元素 cerium（铈）即以 Ceres 命名\(^{[32][c]}\)。

Ceres 的旧天文符号今天仍在占星术中使用，为一柄镰刀，⟨⚳⟩\(^{[22][34]}\)。镰刀是女神 Ceres 的经典象征之一，该符号显然分别由 von Zach 与 Bode 于 1802 年独立提出\(^{[35]}\)。它的形状与 Venus 的符号相似（⟨♀⟩，即一个圆圈下加小十字），但在圆圈处有一处断开。其表现形式存在多种细微变体，包括一种反向图形，以字母 “C”（即 Ceres 的首字母）配以一个加号进行排版。编号圆盘形式的小行星通用符号 ① 于 1867 年引入，并迅速成为惯例\(^{[22][36]}\)。
\subsubsection{分类}
\begin{figure}[ht]
\centering
\includegraphics[width=6cm]{./figures/f71c683b033e055f.png}
\caption{按比例显示的 Ceres（左下）、Moon 与 Earth} \label{fig_Ceres_2}
\end{figure}
\begin{figure}[ht]
\centering
\includegraphics[width=6cm]{./figures/7ea9a2828e334e64.png}
\caption{小行星带中四颗最大小行星的平均直径对比（左侧为矮行星 Ceres）} \label{fig_Ceres_3}
\end{figure}
Ceres 的分类经历过多次变化，并引发过不少争议。Bode 认为 Ceres 是他所提出的位于 Mars 与 Jupiter 之间“失踪行星”\(^{[19]}\)。Ceres 被赋予了行星符号，并在天文学书籍与星表中（与 Pallas、Juno 和 Vesta 一起）作为行星列出，持续了半个多世纪\(^{[37]}\)。

随着更多天体在 Ceres 附近被发现，天文学家开始怀疑它代表了一类新的天体\(^{[19]}\)。当 Pallas 于 1802 年被发现时，Herschel 引入了 “asteroid”（“类星状”）这一术语用于这些天体\(^{[37]}\)，写道：“它们与小星星极为相似，即便使用非常好的望远镜也几乎无法区别”\(^{[38]}\)。1852 年，Johann Franz Encke 在《Berliner Astronomisches Jahrbuch》中宣布，为这些新天体授予传统的行星符号体系过于繁琐，因而引入依照发现顺序在名称前添加数字的新方法。编号体系最初从第五颗小行星 5 Astraea 作为 1 开始，但在 1867 年，Ceres 被纳入该体系，命名为 1 Ceres\(^{[37]}\)。

到 1860 年代，天文学家普遍接受了行星与诸如 Ceres 之类小行星之间存在根本区别，尽管当时 “planet（行星）” 一词尚未被精确定义\(^{[37]}\)。在 1950 年代，科学家普遍停止将大多数小行星视为行星，但 Ceres 有时仍保留这一地位，因为其在地球物理复杂性方面具有类行星特征\(^{[39]}\)。随后在 2006 年，围绕 Pluto 的争议引发了对“行星”定义的呼声，并提出重新分类 Ceres，甚至可能恢复其作为行星的普遍地位\(^{[40]}\)。一项提交给国际天文学联合会（IAU，负责天文命名与分类的全球组织）的提案，将行星定义为：“（a）具有足够质量，其自身引力能克服刚体力使其呈流体静力平衡（近似球形）形状，并且（b）围绕恒星运行，且既非恒星亦非行星的卫星”\(^{[41]}\)。若该决议被采纳，则 Ceres 将成为距 Sun 第五颗行星\(^{[42]}\)，但 2006 年 8 月 24 日大会增加了额外要求，即行星必须“清除其轨道附近的区域”。Ceres 并非行星，因为它无法主导其轨道，与小行星带中成千上万的其他小行星共享轨道，其质量仅占小行星带总质量的约 40\%\(^{[43]}\)。那些满足最初定义但不满足第二条要求的天体（如 Ceres），被分类为矮行星\(^{[44]}\)。行星地质学家仍常常无视这一定义，并依然将 Ceres 视为行星\(^{[45]}\)。

Ceres 被指定为矮行星与小行星。NASA 一篇网页写道，小行星带中第二大天体 Vesta 是最大的小行星\(^{[46]}\)。IAU 在这一问题上态度模糊\(^{[47][48][\text{failed verification}]}\)，但其负责编目相关天体的 Minor Planet Center 指出，矮行星可能具有双重指定\(^{[49]}\)，而 IAU / USGS / NASA 联合地名索引将 Ceres 同时归类为小行星与矮行星\(^{[50]}\)。
\begin{figure}[ht]
\centering
\includegraphics[width=6cm]{./figures/df7ba81daeff6318.png}
\caption{} \label{fig_Ceres_4}
\end{figure}
\subsection{轨道}
\begin{figure}[ht]
\centering
\includegraphics[width=8cm]{./figures/0f8d0cf9a4498022.png}
\caption{Ceres（红色，倾斜）的轨道以及 Jupiter 和内行星（白色与灰色）的轨道。上方示意图从俯视角展示了 Ceres 的轨道。下方示意图为侧视图，显示 Ceres 轨道相对于黄道面的倾角。浅色表示位于黄道面之上；深色表示位于黄道面之下。} \label{fig_Ceres_5}
\end{figure}
Ceres 在 Mars 与 Jupiter 之间运行，位于小行星带中部，其轨道周期为 4.6 个 Earth 年\(^{[2]}\)。与其他行星和矮行星相比，Ceres 的轨道相对 Earth 有适度倾斜；其轨道倾角（\(i\)）为 \(10.6^\circ\)，相比之下 Mercury 为 \(7^\circ\)，Pluto 为 \(17^\circ\)。其轨道也略微拉长，偏心率（\(e\)）为 0.08，而 Mars 为 0.09\(^{[2]}\)。

Ceres 不属于任何小行星族，可能是因为其内部拥有较高比例的冰，具有相同成分的较小天体在太阳系年龄范围内会由于升华而消失殆尽\(^{[51]}\)。曾有人认为 Ceres 属于 Gefion 小行星族\(^{[52]}\)，该族成员具有相似的本征轨道要素，暗示它们可能源于过去一次小行星碰撞。然而，后来发现 Ceres 与 Gefion 家族在成分上不同\(^{[52]}\)，似乎属于“闯入者”，即轨道参数相似但并无共同起源\(^{[53]}\)。
\subsubsection{共振}  
由于小行星带天体质量较小且彼此距离较大，它们之间很少出现引力共振\(^{[54]}\)。尽管如此，Ceres 能够将其他小行星暂时捕获进入 1:1 共振轨道（使其成为暂时的特洛伊天体），持续时间从几十万年至两百多万年不等。目前已识别出 50 个这样的天体\(^{[55]}\)。Ceres 与 Pallas 之间接近 1:1 平均运动轨道共振（其本征轨道周期相差 0.2\%），但不足以在天文学时间尺度上产生显著影响\(^{[56]}\)。
\subsection{自转与轴倾角}
Ceres 的自转周期（即 Ceres 一日）为 9 小时 4 分钟\(^{[10]}\)；其赤道上的小型撞击坑 Kait 被选定为本初子午线\(^{[57]}\)。Ceres 的自转轴倾角为 \(4^\circ\)\(^{[10]}\)，这一数值足够小，使其极区包含永远处于阴影中的撞击坑，这些撞击坑预计会充当冷阱，随时间积累水冰，类似于 Moon 和 Mercury 上发生的情况。预计约 0.14\% 从表面释放的水分子最终会被捕获在这些冷阱中，它们平均“跳跃”三次才会逃逸或被捕获\(^{[10]}\)。

首个进入 Ceres 轨道的探测器 Dawn 确定，Ceres 的北极轴指向赤经 19 h 25 m 40.3 s（291.418°）、赤纬 +66° 45′ 50″（距离 Delta Draconis 约 1.5°），这意味着自转轴倾角为 \(4^\circ\)。这表明 Ceres 目前几乎不存在随纬度变化的季节性阳光差异\(^{[58]}\)。在过去三百万年中，Jupiter 与 Saturn 的引力影响引发了 Ceres 自转轴倾角的周期性变化，范围从 \(2^\circ\) 到 \(20^\circ\)，这意味着过去曾存在季节性日照变化，最近一次出现季节活动大约发生在 14,000 年前。在自转轴最大倾角时期仍能保持阴影的那些撞击坑最有可能在太阳系年龄尺度内保留由喷发或彗星撞击产生的水冰\(^{[59]}\)。
\subsection{地质}
\begin{figure}[ht]
\centering
\includegraphics[width=8cm]{./figures/23fb842f59355847.png}
\caption{谷神星在内太阳系 行星质量物体,按它们从太阳向外的轨道顺序排列 (从左起:水银,维纳斯,地球,the月亮,火星和谷神星)} \label{fig_Ceres_6}
\end{figure}
\begin{figure}[ht]
\centering
\includegraphics[width=8cm]{./figures/18dd605eb8e54f28.png}
\caption{Ceres渲染与搅拌机,从色基图和高度图从NASA和USGS。} \label{fig_Ceres_7}
\end{figure}
谷神星是小行星主带中最大的天体$^{[16]}$。它被分类为 C 型或碳质小行星$^{[16]}$，并且由于存在粘土矿物，还被归类为 G 型小行星$^{[60]}$。它的成分与碳质球粒陨石类似，但并不完全相同$^{[61]}$。它是一个扁球体，其赤道直径比极地直径大 8\%$^{[2]}$。“黎明号”探测器的测量发现，谷神星的平均直径为 939.4 公里（583.7 英里）$^{[2]}$，质量为 $9.38\times10^{20}$ 千克$^{[62]}$。这使得谷神星的密度为 2.16 克/立方厘米$^{[2]}$，表明其约四分之一的质量为水冰$^{[63]}$。

谷神星占据小行星带估计总质量 $(2394\pm5)\times10^{18}$ 千克的 40\%，并且其质量是下一颗小行星——灶神星的 3 又 1/2 倍，但仅为月球质量的 1/78，其表面重力为地球的 1/35（约为月球的 1/6）。它非常接近处于流体静力平衡状态，但仍有一些偏离平衡形状的现象尚未得到解释$^{[64]}$。谷神星是唯一公认轨道周期短于海王星的矮行星$^{[63]}$。建模显示，谷神星的岩石物质可能发生了部分分异，并且可能拥有一个小型核心$^{[65]}$,$^{[66]}$，但数据也与一个由水合硅酸盐构成的地幔且没有核心的结构相一致$^{[64]}$。由于“黎明号”没有配备磁强计，因此目前尚不清楚谷神星是否具有磁场；普遍认为其不具备磁场$^{[67]}$,$^{[68]}$。谷神星内部的分异可能与其缺乏天然卫星有关，因为主带小行星的卫星通常被认为是由碰撞破坏形成，从而产生未分化的碎石堆结构$^{[69]}$。

\subsubsection{表面}
\textbf{成分}
谷神星的表面成分在全球尺度上是均一的，并且富含碳酸盐和经水改变的含铵层状硅酸盐$^{[64]}$，尽管风化层中的水冰含量在极地纬度约为 10\%，而在赤道地区则要干燥得多，甚至几乎不含冰$^{[64]}$。

利用哈勃空间望远镜的研究显示，在谷神星表面存在石墨、硫和二氧化硫。石墨显然是谷神星较古老表面经历空间风化的结果；而后两者在谷神星条件下是挥发性的，理论上应当迅速逃逸或沉积于冷阱中，因此它们显然与相对近期的地质活动有关$^{[70]}$。

在厄努特陨石坑检测到有机化合物$^{[71]}$，并且另有至少 11 个区域被认为可能存在有机化合物$^{[72]}$。该行星近表层大部分区域富含碳，质量分数约为 20\%$^{[73]}$。其碳含量比在地球上分析的碳质球粒陨石高出五倍以上$^{[73]}$。表面碳显示出与岩石—水相互作用产物（如粘土）混合的证据$^{[73]}$。这种化学性质表明谷神星形成于寒冷环境中，可能位于木星轨道之外，并且在水的存在下由超高含碳物质聚合而成，从而可能提供有利于有机化学的条件$^{[73]}$。
\begin{figure}[ht]
\centering
\includegraphics[width=8cm]{./figures/97253d100f841e67.png}
\caption{黑白的谷神星摄影地图，以 180° 经线为中心，并标注官方命名（2017 年 9 月）。} \label{fig_Ceres_8}
\end{figure}
\begin{figure}[ht]
\centering
\includegraphics[width=8cm]{./figures/c0d68798ceb10be6.png}
\caption{谷神星，极地区域（2015 年 11 月）：北极（左）；南极（右）。南极处于阴影中。“Ysolo 山”随后被更名为 “Yamor 山”$^{[74]}$。} \label{fig_Ceres_10}
\end{figure}
\textbf{陨石坑}
\begin{figure}[ht]
\centering
\includegraphics[width=8cm]{./figures/8310f7c97e44cd51.png}
\caption{谷神星的地形图。最低的陨石坑坑底（靛蓝色）与最高的山峰（白色）之间的高程差为 15 公里（10 英里）$^{[75]}$。“Ysolo 山”已被更名为 “Yamor 山”$^{[74]}$。} \label{fig_Ceres_9}
\end{figure}
“黎明号”探测器揭示，谷神星表面存在大量陨石坑，尽管大型陨石坑的数量比预期要少$^{[76]}$。基于当前小行星带形成模型的预测认为，谷神星应当有 10–15 个直径超过 400 公里（250 英里）的陨石坑$^{[76]}$。谷神星上已确认的最大陨石坑为 Kerwan 盆地，直径为 284 公里（176 英里）$^{[77]}$。最可能的原因是地壳发生黏性松弛，使得较大的撞击坑随着时间缓慢变平$^{[76][78]}$。

谷神星的北极区域显示出比赤道区域更多的陨石坑，尤其是东部赤道区域的撞击坑相对较少$^{[79]}$。直径在 20–100 公里（10–60 英里）之间的陨石坑总体数量分布与它们形成于晚期重轰炸时期一致，古老极地区域之外的陨石坑很可能在早期低温火山活动中被抹除$^{[79]}$。三个大型浅盆地（大平原），其边缘已退化，很可能是遭受侵蚀的陨石坑$^{[64]}$。其中最大的 Vendimia 大平原直径达 800 公里（500 英里）$^{[76]}$，同时也是谷神星上最大的单一地理特征$^{[80]}$。这三者之中有两个区域的铵浓度高于平均水平$^{[64]}$。

“黎明号”在谷神星表面观测到 4,423 块直径大于 105 米（344 英尺）的巨石。这些巨石很可能由撞击形成，并出现在陨石坑内或附近，但并非所有陨石坑都包含巨石。大型巨石在高纬度地区数量更多。谷神星上的巨石较为脆弱，并会由于热应力（在黎明和黄昏时，表面温度迅速变化）及陨石撞击而快速降解。它们的最大年龄估计约为 1.5 亿年，远短于灶神星上巨石的存续时间$^{[81]}$。

\textbf{构造特征}

尽管谷神星缺乏板块构造$^{[82]}$，其绝大多数表面特征要么与撞击有关，要么与低温火山活动有关$^{[83]}$，但其表面仍初步识别出一些潜在的构造特征，尤其是在其东半球。Samhain 坑链为谷神星表面尺度达数公里的线状裂隙，它们与撞击没有明显联系，反而更类似于坑链结构，暗示存在埋藏的正断层。此外，谷神星上若干陨石坑的浅层裂纹坑底与低温岩浆侵入相一致$^{[84]}$。

\textbf{低温火山活动}
\begin{figure}[ht]
\centering
\includegraphics[width=6cm]{./figures/998702449cec5764.png}
\caption{阿胡纳山（中央）在最陡的一侧估计有 5 公里（3 英里）高$^{[85]}$。} \label{fig_Ceres_11}
\end{figure}
\begin{figure}[ht]
\centering
\includegraphics[width=6cm]{./figures/7c3452ad060f1bc2.png}
\caption{谷神星上 Cerealia 亮斑与 Vinalia 亮斑的模拟视图} \label{fig_Ceres_12}
\end{figure}
谷神星有一座显著的山峰——阿胡纳山；该山似乎是一座低温火山，并且表面陨石坑很少，表明其最大年龄约为 2.4 亿年$^{[86]}$。其相对较高的重力场表明它密度较大，因此由岩石而非冰构成，而且其位置很可能是由于由盐水与硅酸盐颗粒组成的浆状物在地幔顶部发生底侵作用所致$^{[51]}$。它与 Kerwan 盆地大致呈对跖位置。形成 Kerwan 的撞击所产生的地震能量可能集中在谷神星的另一侧，破裂地壳外层，并触发高黏度低温岩浆（由含盐的泥状冰软化而成）向表面迁移$^{[87]}$。Kerwan 也显示出因撞击导致的地下冰融化而出现液态水作用的证据$^{[77]}$。

一项 2018 年的计算机模拟显示，谷神星上的低温火山一旦形成，会因黏性松弛而在数亿年时间内逐渐消退。研究团队在谷神星表面识别出 22 个强有力的候选特征，可能是经历了松弛作用的低温火山$^{[86][88]}$。Yamor 山是一座古老的、带有撞击陨石坑的山峰，尽管年龄更久远，但由于位于谷神星北极区域，低温阻止了地壳的黏性松弛，因此其外观与阿胡纳山相似$^{[83]}$。模型显示，在过去十亿年中，谷神星平均每 5,000 万年形成一座低温火山$^{[83]}$。火山喷发可能与古老的撞击盆地有关，但在谷神星上分布并不均匀$^{[83]}$。该模型表明，与阿胡纳山的发现相反，谷神星的低温火山必须由密度远低于谷神星地壳平均值的物质构成，否则无法出现已观察到的黏性松弛现象$^{[86]}$。

大量谷神星陨石坑意外地拥有中心坑，这可能与低温火山过程有关；其他陨石坑则有中央峰$^{[89]}$。“黎明号”探测器观察到数百个亮斑（faculae），其中最亮的一处位于直径 80 公里（50 英里）的 Occator 陨石坑中央$^{[90]}$。Occator 中心的亮斑被命名为 Cerealia 亮斑$^{[91]}$；其东侧的亮斑群被命名为 Vinalia 亮斑$^{[92]}$。Occator 有一个宽 9–10 公里的坑，部分被中央穹丘填充。该穹丘形成于亮斑之后，很可能是地下储层结冰导致的，类似于地球北极地区的冰丘（pingos）$^{[93][94]}$。在 Cerealia 上空周期性出现薄雾，这支持了亮斑由某种气体逸出或冰升华所形成的假说$^{[95]}$。2016 年 3 月，“黎明号”在 Oxo 陨石坑首次发现了谷神星表面存在水冰的确凿证据$^{[96]}$。

2015 年 12 月 9 日，NASA 科学家报告称，谷神星亮斑可能由一种来自蒸发盐水的盐类构成，包含六水合硫酸镁（MgSO(_4\cdot6)H(_2)O）；这些亮斑也与富含氨的粘土有关$^{[97]}$。2017 年的近红外光谱分析显示，这些亮斑与大量碳酸钠（Na(_2)CO(_3)）以及少量氯化铵（NH(_4)Cl）或碳酸氢铵（NH(_4)HCO(_3)）相一致$^{[98][99]}$。这些物质被认为起源于到达地表的盐水结晶$^{[100]}$。2020 年 8 月，NASA 确认谷神星是一个富含水的天体，具有深层盐水储库，这些盐水在数百处渗出地表$^{[101]}$，形成“亮斑”，包括出现在 Occator 陨石坑中的那些$^{[102]}$。
\subsubsection{内部结构}
\begin{figure}[ht]
\centering
\includegraphics[width=6cm]{./figures/8a05aad22b023501.png}
\caption{谷神星内部结构的三层模型：外层厚壳（冰、盐类、水合矿物）,富含盐分的液体（盐水）与岩石,“地幔”（水合岩石）} \label{fig_Ceres_13}
\end{figure}
谷神星的活跃地质活动由冰和盐水驱动。从岩石中淋滤出的水估计含有约 5\% 的盐分。总体而言，谷神星按体积计约有 50\% 为水（相比之下地球为 0.1\%），按质量计约有 73\% 为岩石$^{[14]}$。

谷神星最大的一些陨石坑深达数公里，这与富含冰的浅层地下结构不相符。其表面能够保存直径接近 300 公里（200 英里）的陨石坑这一事实表明，谷神星最外层约比水冰强 1000 倍。这与一种由硅酸盐、水合盐和甲烷水合物混合物构成的模型一致，其中水冰体积比例不超过 30\%$^{[64][103]}$。

“黎明号”探测器的重力测量提出了谷神星内部结构的三种互相竞争的模型$^{[14]}$。在三层模型中，谷神星被认为由一个约 40 公里（25 英里）厚的外层地壳组成，其成分包括冰、盐类和水合矿物；内部为泥状的“地幔”，由水合岩石（例如粘土）构成，两者之间隔着一层约 60 公里（37 英里）厚的泥状盐水与岩石混合层$^{[104]}$。目前无法确定谷神星深部内部是否含有液态物质，或是否存在一个富含金属的高密度核心$^{[104]}$，但其较低的中心密度表明其内部可能保留约 10\% 的孔隙度$^{[14]}$。一项研究估计，核层与地幔/地壳层的密度分别为 2.46–2.90 与 1.68–1.95 g/cm³，而地幔与地壳合计厚度为 70–190 公里（40–120 英里）。预计核心只会发生部分脱水（冰的排出），但地幔相对于水冰的高密度反映其富含硅酸盐和盐类$^{[9]}$。也就是说，核心（若存在）、地幔和地壳均由岩石和冰组成，只是比例不同。

谷神星的矿物组成只能（间接地）确定其外层约 100 公里（60 英里）的区域。其固体外壳厚约 40 公里（25 英里），为冰、盐类和水合矿物的混合物。其下可能存在一薄层含有少量盐水，该结构延伸至至少 100 公里（60 英里）的探测深度以下，再往下被认为是主要由水合岩石（如粘土）构成的地幔$^{[104]}$。

在一种两层模型中，谷神星由一个球粒状物质组成的核心和一个由冰与微米级固体微粒（即“泥”）构成的地幔组成。表面冰的升华将留下约 20 米厚的水合微粒沉积层。区分程度的范围与数据相符：从一个大型核心（直径 360 公里〔220 英里〕，其中 75\% 为球粒、25\% 为微粒，地幔成分为 75\% 冰、25\% 微粒），到一个小型核心（直径 85 公里〔55 英里〕，几乎完全由微粒组成，地幔为 30\% 冰、70\% 微粒）。若核心较大，则核心—地幔边界温度应足够高，可形成盐水囊；若核心较小，则在 110 公里（68 英里）以下地幔应保持为液态。在后者情况下，液体储层中 2\% 的冻结就足以产生压力迫使部分液体上升至表面，造成低温火山活动$^{[105]}$。

另一种两层模型认为谷神星发生了部分分异，即形成一个富含挥发物的地壳与一个较致密的水合硅酸盐地幔。根据可能撞击过谷神星的陨石类型，可以计算一系列可能的地壳与地幔密度。若假定为 CI 类陨石（密度 2.46 g/cm³），则地壳厚度约 70 公里（40 英里），密度为 1.68 g/cm³；若为 CM 类陨石（密度 2.9 g/cm³），则地壳厚度约 190 公里（120 英里），密度为 1.9 g/cm³。拟合最佳的模型得出地壳厚度约 40 公里（25 英里），密度约 1.25 g/cm³，而地幔/核心的密度约为 2.4 g/cm³$^{[64]}$。
\subsection{外逸层}
2017 年，“黎明号”确认谷神星存在短暂的水蒸气大气$^{[106]}$。早在 2014 年初就已出现大气迹象，当时赫歇尔空间天文台在谷神星上探测到局部的中纬度水蒸气来源，其范围不超过 60 公里（40 英里），每秒释放约 (10^{26}) 个（3 千克）水分子$^{[107][108]}$。两个潜在来源区域——Piazzi 区（东经 123°，北纬 21°）与 A 区（东经 231°，北纬 23°）——在 W. M. Keck 天文台的近红外观测中显示为暗区（A 区具有明亮中心）。可能的蒸汽释放机制包括约 0.6 平方公里（0.2 平方英里）裸露表面冰的升华、由放射性内部热驱动的低温火山喷发$^{[107]}$，或由于上层冰变厚导致地下海洋受压$^{[111]}$。2015 年，David C. Jewitt 将谷神星列入活跃小行星名单$^{[112]}$。在距太阳小于 5 AU 的位置，表面冰处于不稳定状态$^{[113]}$，因此直接暴露在太阳辐射下预期会升华。太阳耀斑和日冕物质抛射所产生的质子会溅射表面裸露冰，使得水蒸气探测与太阳活动呈正相关$^{[114]}$。水冰可以从谷神星深部迁移至表面，但会在短时间内逃逸。当谷神星远离太阳时，表面升华预期会降低，而内部驱动的释放则不受轨道位置影响。先前有限的数据暗示谷神星升华类似彗星$^{[107]}$，但“黎明号”的证据表明地质活动至少部分负责这一现象$^{[115]}$。

利用“黎明号”的伽马射线与中子探测器（GRaND）进行的研究显示，谷神星能加速来自太阳风的电子；最被接受的假说认为，这些电子通过太阳风与稀薄水蒸气外逸层的碰撞被加速$^{[116][117]}$。这种弓形激波也可能由一个短暂的磁场解释，但这种可能性较小，因为谷神星内部被认为没有足够的电导率$^{[117]}$。谷神星的稀薄外逸层通过撞击暴露的冰斑、冰通过多孔冰壳的扩散以及太阳活动期间质子溅射不断得到补充$^{[118][119][120]}$。蒸汽扩散速率随颗粒大小变化$^{[121]}$，并受一个由约 1 微米粒子聚集而成的全球尘埃覆盖层的强烈影响$^{[122]}$。仅靠升华实现的外逸层补给非常微小，目前的气体释放率仅为 0.003 千克/秒$^{[119]}$。有关外逸层存在的多种模型已被提出，包括弹道轨迹模型、DSMC 模型与极地帽数值模型$^{[123][120][124]}$。结果显示，弹道轨迹模型得到的水外逸层半衰期为 7 小时；DSMC 模型在光学稀薄大气下得到了持续数十天且释放率为 6 千克/秒的情形；极地帽模型显示，由外逸层输送的水会形成季节性极地冰盖。外逸层中水分子的运动主要受弹跳式运动与表面相互作用的支配，但对其与行星风化层直接相互作用的了解较少$^{[119]}$。
\subsection{起源与演化}
谷神星是一颗幸存的原行星，形成于 45.6 亿年前；与智神星（Pallas）与灶神星一起，是太阳系内层中仅存的三颗原行星之一$^{[125]}$，其余要么在碰撞中合并形成类地行星，要么在撞击中被破碎$^{[126]}$，或被木星抛出轨道$^{[127]}$。尽管谷神星当前的位置位于小行星带，但其成分并不符合在该区域形成的特征。相反，谷神星似乎形成于木星与土星轨道之间，并在木星向外迁移时被偏转进入小行星带$^{[14]}$。在 Occator 陨石坑中发现的铵盐支持这一源自外太阳系的假说，因为氨在外太阳系丰富得多$^{[128]}$。

谷神星早期地质演化依赖于其形成期间与形成后的热源：来自星子聚积的撞击能量以及放射性核素（可能包括已灭绝的短寿命核素，如 (^{26})Al）的衰变。这些热源可能足以使谷神星分化形成岩质行星核与冰质地幔，甚至形成液态海洋$^{[64]}$，这一过程在其形成后不久就可能发生$^{[66]}$。这一海洋在冻结时应在表面之下留下冰层。“黎明号”没有发现这类证据，表明谷神星原始地壳至少部分被后期撞击破坏，使得冰与盐类以及富含硅酸盐的古代海底物质深度混合$^{[64]}$。

谷神星拥有令人惊讶的少量大型陨石坑，表明黏性松弛与低温火山活动已抹去了较古老的地质特征$^{[129]}$。粘土与碳酸盐的存在需要高于 50°C 的化学反应，符合热液活动条件$^{[51]}$。

谷神星的地质活动随时间显著减弱，其表面主要由撞击坑主导；然而，“黎明号”的证据显示，内部过程仍持续在相当程度上塑造谷神星表面$^{[130]}$，这与此前预测谷神星因其较小体积而在早期即终止内部地质活动的观点相反$^{[131]}$。
\subsection{可居性}
\begin{figure}[ht]
\centering
\includegraphics[width=6cm]{./figures/7721eaf63a85b575.png}
\caption{风化层上部一米范围内的氢浓度（蓝色）显示存在水冰} \label{fig_Ceres_14}
\end{figure}
尽管谷神星并不像火星、木卫二（Europa）、土卫二（Enceladus）或土卫六（Titan）那样常被积极讨论为潜在的地外微生物生命栖息地，但在地球之后，它是太阳系内层中含水量最多的天体$^{[51]}$，其地下可能存在的盐水囊能够为生命提供栖息环境$^{[51]}$。与木卫二或土卫二不同，谷神星不经历潮汐加热，但它距离太阳足够近，并含有数量足够的长寿命放射性同位素，可以在其地下长时间维持液态水$^{[51]}$。对有机化合物的遥感探测，以及其近表层水与 20\% 质量分数碳混合的存在，可能提供适合有机化学的条件$^{[73]}$。在生化元素中，谷神星富含碳、氢、氧和氮$^{[132]}$，但尚未探测到磷$^{[133]}$，而硫虽在哈勃紫外观测中有所提示，却未被“黎明号”探测到$^{[51]}$。
\subsection{观测与探测}
\subsubsection{观测}
\begin{figure}[ht]
\centering
\includegraphics[width=6cm]{./figures/b8ad9e533ff8f755.png}
\caption{} \label{fig_Ceres_15}
\end{figure}
当位于近日点附近的冲位置时，谷神星的视星等可以达到 +6.7$^{[134]}$。这对于一般裸眼来说仍然太暗而难以看见，但在理想观测条件下，视力敏锐者可能可以看到它。灶神星是唯一一颗能够定期达到类似亮度视星等的小行星，而智神星与 7 号鸢星只有在同时位于冲位置且接近日点时才能达到这样的亮度$^{[135]}$。位于合位置时，谷神星的视星等约为 +9.3，这相当于使用 10×50 双筒望远镜所能观测到的最暗天体；因此，在新月前后自然黑暗且晴朗的夜空下，使用这样的望远镜可以看到谷神星$^{[17]}$。

1984 年 11 月 13 日，谷神星掩食恒星 BD+8°471 的现象在墨西哥、佛罗里达及整个加勒比地区被观测到，从而获得了关于谷神星大小、形状与反照率的更优测量$^{[136]}$。1995 年 6 月 25 日，哈勃以 50 公里（30 英里）分辨率获得了谷神星的紫外图像$^{[60]}$。2002 年，W. M. Keck 天文台借助自适应光学技术，以 30 公里（20 英里）分辨率获得了红外图像$^{[137]}$。

在“黎明号”任务之前，只有极少数谷神星表面特征得到明确探测。1995 年的高分辨率紫外哈勃图像显示其表面存在一处暗斑，被昵称为“Piazzi”，以纪念谷神星的发现者$^{[60]}$。被认为是一处陨石坑。2003 年和 2004 年哈勃拍摄的谷神星一整圈自转的可见光图像显示了 11 处可识别的表面特征，但其性质不明$^{[13][138]}$。其中一处与“Piazzi”特征相对应$^{[13]}$。2012 年，Keck 天文台使用自适应光学拍摄的谷神星整圈自转近红外图像显示，亮斑与暗斑随谷神星自转而移动$^{[139]}$。其中两处暗斑呈圆形，被认为是陨石坑；其中一处有明亮中心，另一处被认定为“Piazzi”特征$^{[139]}$。“黎明号”最终揭示，“Piazzi”是位于 Vendimia 大平原中央、靠近 Dantu 陨石坑的一片暗区，而另一处暗斑位于 Hanami 大平原内，靠近 Occator 陨石坑$^{[140]}$。
\subsubsection{“黎明号”任务}
\begin{figure}[ht]
\centering
\includegraphics[width=6cm]{./figures/b6cab2dfce41843e.png}
\caption{} \label{fig_Ceres_17}
\end{figure}
\begin{figure}[ht]
\centering
\includegraphics[width=6cm]{./figures/390426e27635919c.png}
\caption{艺术家的概念黎明航天器} \label{fig_Ceres_16}
\end{figure}
20 世纪 90 年代初，NASA 启动了“探索者计划”，该计划旨在开展一系列低成本的科学任务。1996 年，该计划的研究团队提出了一项高优先级任务，使用带有离子推进器的航天器探索小行星带。资金问题困扰该计划近十年之久，但到 2004 年，“黎明号”飞行器通过了关键设计评审$^{[141]}$。

“黎明号”是首个访问灶神星或谷神星的太空任务，于 2007 年 9 月 27 日发射。2011 年 5 月 3 日，“黎明号”在距灶神星 1,200,000 公里（750,000 英里）处获得其首张目标图像$^{[142]}$。在围绕灶神星运行 13 个月后，“黎明号”利用离子推进器启程飞往谷神星，并于 2015 年 3 月 6 日以 61,000 公里（38,000 英里）的距离被谷神星引力捕获$^{[143][144]}$，这比“新视野号”飞越冥王星早了四个月$^{[144]}$。

该航天器的仪器包括测绘相机、可见光与红外光谱仪，以及伽马射线与中子探测器。这些仪器用于研究谷神星的形状与元素组成$^{[145]}$。2015 年 1 月 13 日，随着“黎明号”接近谷神星，该探测器拍摄了其第一批接近哈勃分辨率的图像，显示了撞击坑和表面的一处高反照率小亮斑。从 2 月至 4 月，开展了多轮成像，每次分辨率逐步提高$^{[146]}$。

“黎明号”的任务方案要求它从一系列连续降低高度的圆形极轨道对谷神星进行研究。2015 年 4 月 23 日，它进入了距谷神星 13,500 公里（8,400 英里）的第一观测轨道（“RC3”），只停留一个轨道周期（15 天）$^{[147][148]}$。随后航天器降低轨道至 4,400 公里（2,700 英里），用于第二观测轨道（“survey”）为期三周$^{[149]}$，然后降低至 1,470 公里（910 英里）（“HAMO”，高轨测绘轨道）为期两个月$^{[150]}$，再进一步降低至最终轨道 375 公里（233 英里）（“LAMO”，低轨测绘轨道），停留至少三个月$^{[151]}$。2015 年 10 月，NASA 发布了“黎明号”拍摄的谷神星真彩色肖像图$^{[152]}$。2017 年，“黎明号”任务获延长，执行一系列更接近谷神星的轨道，直到用于维持轨道的肼燃料耗尽$^{[153]}$。

“黎明号”很快发现了低温火山活动的证据。2015 年 2 月 19 日的图像中，在一处陨石坑内看到两处独立的亮斑（或高反照率特征，与早期哈勃图像中看到的亮斑不同）$^{[154]}$，引发了关于可能的低温火山成因$^{[155]}$或气体逸出的猜测$^{[156]}$。2016 年 9 月 2 日，“黎明号”团队的科学家在《Science》上发表论文，认为阿胡纳山是谷神星上迄今最强有力的低温火山证据$^{[87]}$。2015 年 5 月 11 日，NASA 发布了分辨率更高的图像，显示这些亮斑由多个较小的亮斑组成$^{[157]}$。2015 年 12 月 9 日，NASA 科学家报告称，谷神星的亮斑可能与某种盐类有关，特别是一种含有六水合硫酸镁（MgSO(_4\cdot6)H(_2)O）的盐水形式；这些亮斑也与富含氨的粘土有关$^{[97]}$。2016 年 6 月，根据近红外光谱结果，这些亮斑含有大量碳酸钠（$Na_2 CO_3$），表明近期地质活动可能参与了亮斑的形成$^{[158]}$。

2018 年 6 月至 10 月，“黎明号”在距谷神星 35 公里（22 英里）至 4,000 公里（2,500 英里）之间的轨道上运行$^{[159]}$。“黎明号”任务于 2018 年 11 月 1 日结束，此时航天器燃料耗尽$^{[160]}$。
\subsubsection{未来任务}
2020 年，欧洲航天局团队提出了 Calathus 任务概念，这是一项针对 Occator 陨石坑的后续任务，旨在将其亮色碳酸盐亮斑与暗色有机物样本带回地球$^{[161]}$。中国国家航天局正在设计一项取回谷神星样本的任务，计划在 2020 年代实施$^{[162]}$。
\subsection{参见}
\begin{itemize}
\item 异常小行星列表
\item 按大小排列的太阳系天体列表
\item 前行星列表
\end{itemize}
\subsection{注释}
a.根据已知参数计算：
\begin{itemize}
\item 表面积：($4\pi r^{2}$)
\item 表面重力：($\frac{GM}{r^{2}}$)
\item 逃逸速度：($\sqrt{\frac{2GM}{r}}$)
\item 自转速度：($\frac{\text{周长}}{\text{自转周期}}$)
\end{itemize}\\
b.对于谷神星给出的数值是平均转动惯量，这被认为比极转动惯量更能代表其内部结构，因为谷神星的极向扁平程度很高$^{[9]}$。\\
c.1807 年，克拉普罗特试图将该元素名称改为 cererium，以避免与词根 *cēra* “蜡”（如 *cereous*，“蜡质的”）混淆，但该名称并未流行起来$^{[33]}$。\\
d.与基于潮汐作用产生的土卫二（较小天体）和木卫二（较大天体）喷流相比，谷神星的该逸出速率较为温和，分别为 200 千克/秒$^{[109]}$ 与 7000 千克/秒$^{[110]}$。
\subsection{参考资料}
\begin{enumerate}
\item Schmadel, Lutz(2003)。小行星名称词典(第5版)。德国: 施普林格。第15页。ISBN 978-3-540-00238-3。已存档从原件到2021年2月16日。已检索1月21日2021。
\item “JPL小型数据库浏览器: 1 Ceres”。JPL太阳能系统动力学。已存档从2021年6月9日开始。已检索9月26日2021。
\item “在新的谷神星上”。自然哲学，化学和艺术杂志。1802.已存档从2022年5月29日开始。已检索5月29日2022年。
\item 苏阿米，D.；Souchay, J.(2012年7月)。“太阳系的不变平面”。天文学与天体物理学。543: 11。Bibcode:2012A & A...543A.133S。doi:10.1051/0004-6361/201219011。A133.
\item “AstDyS-2谷神星合成适当轨道元素”。意大利比萨大学数学系。已存档从2011年11月21日的原件。已检索10月1日2011年。
\item Ermakov, A.I.;傅，R。R、；卡斯蒂略-罗盖兹，J.C.;雷蒙德，C。A.;公园，R。美国；Preusker, F.;罗素，C。T.;史密斯，D。E.;Zuber, M.T.(2017年11月)。“对ceres内部结构的限制以及由Dawn航天器测量的形状和重力的演变”。地球物理研究杂志: 行星。122(11):2267-2293。Bibcode:2017JGRE..122.2267E。doi:10.1002/2017JE005302。S2CID133739176。 
\item 公园，R.S.；沃恩，a.t.；Konopliv, A.S.;埃尔马科夫，A.I.；Mastrodemos, N.;Castillo-Rogez, J.C.;乔伊，s.p.；Nathues, A.;波兰斯基，C.A.；雷曼，医学博士; 里德尔，J.E.；雷蒙德，C.A.；罗素，C.T.；祖伯，M.T.(2019年2月)。“使用黎明成像数据的立体测斜术的谷神星高分辨率形状模型”。伊卡洛斯。319:812-827。Bibcode:2019年Icar .. 319 .. 812便士。doi:10.1016/j.icarus.2018.10.024。S2CID126268402。 
\item 毛X。；麦金农，W.B.(2018)。“更快的古spin和深层的无补偿质量可能解释了ceres今天的形状和重力”。伊卡洛斯。299:430-442。Bibcode:2018年Icar .. 299 .. 430米。doi:10.1016/j.icarus.2017.08.033。
\item 公园，R。美国；Konopliv, A.美国；比尔，B.G.;兰博克斯，N.；卡斯蒂略-罗盖兹，J.C.;雷蒙德，C。A.;沃恩，A.T.;Ermakov, A.I.;Zuber, M.T.;傅，R。R、；Toplis, M.J.;罗素，C。T.;Nathues, A.;Preusker，F。(2016年8月3日)。“从其重力场和形状推导的 (1) 谷神星的部分微分内部”。自然。537(7621):515-517。Bibcode:2016 Natur.537..515P。doi:10.1038/自然18955。PMID27487219。S2CID4459985。  
\item Schorghofer, N.;马扎里科，E.；普拉茨，T.；Preusker, F.;施罗德，S.E.;雷蒙德，C。A.;罗素，C。T.(2016年7月6日)。“矮行星谷神星的永久阴影区域”。地球物理研究快报。43(13):6783-6789。Bibcode:2016乔治。。43.6783S。doi:10.1002/2016GL069368。
\item Konopliv, A.S.;公园，R.S.；沃恩，a.t.；Bills, B.G.;阿斯玛，S.W.；埃尔马科夫，A.I.；兰博克斯，N.；雷蒙德，C.A.；Castillo-Rogez, J.C.;罗素，C.T.；史密斯，D.E.；祖伯，M.T.(2018)。“谷神星重力场、自旋极、自转周期和轨道来自黎明的辐射跟踪和光学数据”。伊卡洛斯。299:411-429。Bibcode:2018年Icar .. 299 .. 411k。doi:10.1016/j.icarus.2017.08.005。
\item "小行星Ceres P_constants (PcK) SPICE内核文件"。NASA导航和辅助信息设施。已存档从原件到2020年7月28日。已检索9月8日2019。
\item 李建阳; 麦克法登，露西·A；帕克，乔尔·沃姆。(2006)。“1谷神星的光度分析和HST观测的表面测绘”。伊卡洛斯。182(1):143-160。Bibcode:2006Icar .. 182 .. 143L。doi:10.1016/j.icarus.2005.12.012。
\item 罗杰斯，J.C.卡斯蒂略; 雷蒙德，C。A.;罗素，C。T.;团队，黎明 (2017)。"黎明在Ceres: 我们学到了什么？”(PDF)。NASA, JPL。已存档(PDF)从2018年10月8日的原件。已检索7月19日2021。
\item 托西，F.；卡普里亚，M。T.; 等人 (2015年)。“矮行星谷神星的表面温度: 黎明的初步结果”。第46届月球与行星科学会议: 11960。Bibcode:2015EGUGA..1711960T。已检索5月25日2021。
\item 里夫金，A.美国；Volquardsen, E.L.;克拉克，B。E.(2006)。“谷神星的表面组成: 碳酸盐和富铁粘土的发现”(PDF)。伊卡洛斯。185(2):563-567。Bibcode:2006年Icar。。185。563R。doi:10.1016/j.icarus.2006.08.022。已存档(PDF)从2007年11月28日的原件。已检索12月8日2007。
\item 国王，鲍勃 (2015年8月5日)。“让我们认真对待Ceres”。天空和望远镜。已检索7月25日2022年。
\item “小行星 (1) 谷神星-摘要”。AstDyS-2。已存档从原件到2020年7月26日。已检索10月15日2019。
\item 霍斯金，迈克尔 (1992年6月26日)。“博德定律与谷神星的发现”。巴勒莫天文台 “朱塞佩S.Vaiana ”。存档自原来的2007年11月16日。已检索7月5日2007。
\item 霍格，海伦·索耶 (1948)。“提tius-bode定律与谷神星的发现”。加拿大皇家天文学会杂志。242:241-246。Bibcode:1948JRASC..42..241S。已存档从原件到2021年7月18日。已检索7月18日2021。
\item 伊丽莎白·兰道 (2016年1月26日)。“谷神星: 保守秘密215年”。NASA。已存档从原始的2019年5月24日。已检索1月26日2016。
\item 福布斯，埃里克G。(1971)。“高斯与谷神星的发现”。天文学史杂志。2(3):195-199。Bibcode:1971JHA.....2..195F。doi:10.1177/002182867100200305。S2CID125888612。已存档从原件到2021年7月18日。已检索7月18日2021。 
\item 坎宁安，克利福德J.(2001)。第一颗小行星: 谷神星，1801-2001。明星实验室出版社。ISBN 978-0-9708162-1-4。已存档从2016年5月29日开始。已检索10月23日2015。
\item 迈克尔·马丁·涅托 (1972)。Titius-bode行星距离定律: 其历史和理论。佩加蒙出版社。ISBN 978-1-4831-5936-2。已存档从原件到2021年9月29日。已检索9月23日2021。
\item 休斯，大卫·W (1994)。“前四颗小行星直径的历史解开”。皇家天文学会季刊。35:331-344.Bibcode:1994QJRAS..35..331H。已存档从原始的2021年8月2日。已检索8月2日2021。
\item Foderà Serio, G.;马纳拉，A.；Sicoli, P.(2002)。“朱塞佩·皮亚齐和谷神星的发现”(PDF)。在W。F.Bottke Jr.; A.切利诺; P。Paolicchi; R.P。宾泽尔 (编辑)。小行星III。图森: 亚利桑那大学出版社。pp. 17-24.已存档(PDF)从原件到2012年4月16日。已检索6月25日2009年。
\item Rüpke, Jörg(2011)。罗马宗教的同伴。约翰·威利和儿子。pp. 51-52。ISBN 978-1-4443-4131-7。已存档从2015年11月15日的原件。已检索10月23日2015。
\item “黎明飞船在灶神星上发现了水的痕迹”。科技日报。2012年9月21日。已存档从原件到2021年9月23日。已检索9月23日2021。
\item 里夫金，A.等。(2012)。“谷神星的表面组成”。在罗素，克里斯托弗; 雷蒙德，卡罗尔 (编辑)。黎明任务到小行星4 Vesta和1 Ceres。斯普林格。第109页。ISBN 978-1-4614-4902-7。
\item 威廉·托马斯·桑顿 (2012) [1878]。“Epode 16”。从贺拉斯逐字逐句。纳布出版社。第314页。ISBN 978-1-279-56080-8。
\item Booth, W.(1823)。罗马诗歌的花朵。哈佛大学。
\item “铈: 历史信息”。自适应光学。已存档从2010年4月9日的原件。已检索4月27日2007。
\item "铈"。牛津英语词典(在线版)。牛津大学出版社。 (订阅或参与机构成员必需。)
\item JPL/NASA (2015年4月22日)。“什么是矮行星？”。喷气推进实验室。已存档从原始的2021年12月8日。已检索1月19日2022年。
\item 坎宁安，克利福德 (2015)。第一颗小行星谷神星的发现。斯普林格国际公司第69、164、206页。ISBN 978-3-319-21777-2。OCLC 1100952738。
\item 古尔德，B。A.(1852)。“关于小行星的符号符号”。天文杂志。2(34): 80。Bibcode:1852AJ......2...80G。doi:10.1086/100212。
\item 希尔顿，詹姆斯L.(2001年9月17日)。“小行星是什么时候变成小行星的？”。美国海军天文台。存档自原来的2007年11月6日。已检索8月16日2006。
\item 威廉·赫歇尔(1802年5月6日)。“对两个最近发现的天体的观测”。伦敦皇家学会的哲学交易。92:213-232。Bibcode:1802RSPT...92..213H。doi:10.1098/rstl.1802.0010。JSTOR107120。S2CID115664950。  
\item 梅茨格，菲利普·T.；塞克斯，马克·V。；斯特恩，艾伦; 鲁尼翁，柯比 (2019年)。“小行星从行星到非行星的重新分类”。伊卡洛斯。319:21-32。arXiv:1805.04115。Bibcode:2019Icar .. 319...21米。doi:10.1016/j.icarus.2018.08.026。S2CID119206487。 
\item 康纳，史蒂夫 (2006年8月16日)。“太阳系将迎来三个新行星”。《新西兰先驱报》。已存档从原件到2021年7月19日。已检索7月19日2021。
\item 欧文·金格里奇等 (2006年8月16日)。IAU对 “行星” 和 “Plutons” 的定义草案"。IAU.存档自原来的2008年8月27日。已检索4月27日2007。
\item “IAU对行星和岩体的定义草案”。太空日报。2006年8月16日。已存档从2009年9月6日的原件。已检索4月27日2007。
\item 皮特耶娃，E.V.(2018)。“来自行星和航天器运动的主要小行星带和柯伊伯带的质量”。太阳系研究。44(8-9):554-566.arXiv:1811.05191。Bibcode:2018年AstL...44 .. 554P。doi:10.1134/S1063773718090050。S2CID119404378。 
\item “深度 | Ceres”。NASA太阳系探索。2017年11月9日已存档从原件到2019年4月21日。已检索4月21日2019。
\item 梅茨格，菲利普·T.；格伦迪，W。M、；马克·V·赛克斯；艾伦·斯特恩; 詹姆斯·F·贝尔三世；夏琳·德特里奇；鲁尼翁，柯比; 萨默斯，迈克尔 (2022)。“卫星是行星: 行星科学分类学中的科学有用性与文化目的论”。伊卡洛斯。374114768。arXiv:2110.15285。Bibcode:2022年Icar .. 37414768米。doi:10.1016/j.icarus.2021.114768。S2CID240071005。已检索8月8日2022年。 
\item “一个任务，两个非凡的目的地”。科学。NASA(science.nasa.gov)。2017年10月26日。已存档从原件到2020年7月17日。已检索7月14日2020。小行星的大小从灶神星不等-最大的直径约为329英里 (530公里)。
\item Lang，Kenneth (2011)。剑桥太阳系指南。剑桥大学出版社，第372、442页。ISBN 978-1-139-49417-5。已存档从原件到2020年7月26日。已检索7月27日2019-通过Google.
\item “问题和答案2”。iau.org。国际天文学联合会。已存档从原来的2016年1月30日。已检索1月31日2008。谷神星是 (或者现在我们可以说它是) 最大的小行星。..还有许多其他小行星可以接近谷神星的轨道。
\item Spahr, T.B.(2006年9月7日)。"编辑通知"。minorplanetcenter.net。MPEC 2006-r19。小行星中心(MPC)。已存档从2008年10月10日的原件。已检索1月31日2008。“矮行星” 的编号并不排除它们在此类物体的可能单独目录中具有双重名称。
\item "目标: 谷神星"。行星命名法。planetarynames.wr.usgs.gov。国际天文学联合会/USGS天体地质科学中心/国家航空航天局。已存档从2017年10月13日的原件。已检索9月27日2021。
\item Castillo-rogez，Julie C.; 等人 (2020年1月31日)。“谷神星: 天体生物学目标和可能的海洋世界”。天体生物学。20(2):269-291。Bibcode:2020年阿斯比奥。。20。。269C。doi:10.1089/ast.2018.1999。PMID31904989。 
\item Cellino，A.; 等人 (2002年)。“小行星家族的光谱特性”(PDF)。小行星III。亚利桑那大学出版社。第633-643页 (表在第636页)。Bibcode:2002年aste.book .. 633C。已存档(PDF)从2016年3月28日的原件。已检索8月6日2011年。
\item 凯利，M。美国；加菲，M。J.(1996)。“谷神星 (威廉姆斯67号) 小行星家族的遗传研究”。美国天文学会公报。28: 1097。Bibcode:1996DPS....28.1009K。
\item Christou, A.A.(2000)。“主要小行星带的共轨道物体”。天文学与天体物理学。356:L71-L74。Bibcode:2000A & A...356L..71C。
\item Christou, A.A.;维格特，P。(2012年1月)。“与谷神星和灶神星共同运行的主带小行星种群”。伊卡洛斯。217(1):27-42。arXiv:1110.4810。Bibcode:2012Icar .. 217...27C。doi:10.1016/j.icarus.2011.10.016。ISSN0019-1035。S2CID59474402。  
\item Kovačević, A.B.(2011)。“根据最有效的引力近距离接触确定谷神星的质量”。皇家天文学会月刊。419(3):2725-2736。arXiv:1109.6455。Bibcode:2012MNRAS.419.2725K。doi:10.1111/j.1365-2966.2011.19919.x。
\item 雷曼，马克 (2015年10月30日)。谷神星的新地图揭示了神秘 “亮点” 周围的地形"。NASA。已检索9月13日2022年。
\item 罗素，C。T.;雷蒙德，C。A.; 等人 (2015年7月21日)。“05。黎明从调查轨道上探索谷神星的结果“(PDF)。NASA.已存档(PDF)从原来的2015年9月5日。已检索9月23日2021。
\item “Ceres的阴影陨石坑中的冰与倾斜历史有关”。NASA太阳系探索。2017。已存档从原件到2021年5月15日。已检索5月15日2021。
\item 帕克，J.W.;艾伦·S·斯特恩；托马斯·彼得·C等人 (2002年)。“用哈勃太空望远镜从紫外线观测中分析谷神星的第一批磁盘分辨图像”。天文杂志。123(1):549-557。arXiv:astro-ph/0110258。Bibcode:2002AJ....123..549P。doi:10.1086/338093。S2CID119337148。 
\item 托马斯·B·麦考德；Zambon，Francesca (2019年1月15日)。“黎明任务中谷神星的表面成分”。伊卡洛斯。318:2-13.Bibcode:2019年Icar .. 318。2米。doi:10.1016/j.icarus.2018.03.004。S2CID125115208。已存档从2021年5月20日开始。已检索7月25日2021。 
\item 雷曼，马克·D.(2015年5月28日)。“黎明杂志，2015年5月28日”。喷气推进实验室。存档自原来的2015年5月30日。已检索5月29日2015。
\item Nola Taylor Redd (2018年5月23日)。“谷神星: 最小和最近的矮行星”。space.com。已存档从原件到2021年9月5日。已检索7月25日2021。
\item 雷蒙德，C.；卡斯蒂略-罗盖兹，J.C.;公园，R。美国；Ermakov，A.; 等人 (2018年9月)。黎明数据揭示了谷神星复杂的地壳演化(PDF)。欧洲行星科学大会。卷12.已存档(PDF)从原件到2020年1月30日。已检索7月19日2020。
\item 诺伊曼，W.；布鲁尔，D.；Spohn, T.(2015年12月2日)。“对谷神星的内部结构进行建模: 通过蠕变将吸金与压实相结合，并对水岩分化产生影响”(PDF)。天文学与天体物理学。584: a117。Bibcode:2015A & A...584A.117N。doi:10.1051/0004-6361/201527083。已存档(PDF)从2016年8月22日的原件。已检索7月10日2016。
\item 巴蒂亚，G.K.；Sahijpal, S.(2017)。“早期太阳系中跨海王星天体，冰冷卫星和小型冰冷行星的热演化”。气象学与行星科学。52(12):2470-2490。Bibcode:2017M & p。52.2470b。doi:10.1111/maps.12952。S2CID133957919。 
\item 罗素，C.T.；比利亚雷亚尔，m.n.；Prettyman, T.H.;山下，N。(2018年5月16日)。“太阳风与灶神星和谷神星的相互作用: 对它们磁矩的影响”。欧空局宇宙。已检索10月10日2022年。
\item 诺德海姆，T.A.；Castillo-Rogez, J.C.;比利亚雷亚尔，m.n.；史高丽，J.E.C.；Costello, E.S.(2022年5月1日)。“谷神星的辐射环境及其对表面采样的影响”。天体生物学。22(5):509-519。Bibcode:2022年阿斯比奥。。22 .. 509N。doi:10.1089/ast.2021.0080。ISSN1531-1074。PMID35447049。S2CID248323790。已存档从2022年4月25日开始。已检索7月22日2022年。   
\item 露西·麦克法登；大卫·R·斯基尔曼；Memarsadeghi, N.(2018年12月)。黎明任务寻找谷神星: 完整的原行星没有卫星伊卡洛斯。316:191-204.Bibcode:2018年Icar .. 316 .. 191米。doi:10.1016/j.icarus.2018.02.017。S2CID125181684。 
\item “在谷神星上观察到的硫，二氧化硫，石墨化碳”。spaceref.com。2016年9月3日存档自原来的2021年9月29日。已检索9月8日2016。
\item 汉娜·H·卡普兰; 拉尔夫·E·米利肯；亚历山大，康奈尔M.O'D。(2018年5月21日)。“对谷神星上有机物的丰度和组成的新限制”。地球物理研究快报。45(11):5274-5282。Bibcode:2018乔治。。45.5274k。doi:10.1029/2018GL077913。S2CID51801398。 
\item Rizos, J.L.;阳光，J。M、；戴利，R.T.;Nathues, A.;De Sanctis, C.;拉博尼，A.；帕斯卡特，J.H、; 法纳姆，T。L.;Kloos, J.;奥尔蒂斯，J.L.(2024)。“谷神星上富含有机物的地区的新候选者”。行星科学杂志。5(12): 262。arXiv:2406.14893。Bibcode:2024PSJ.....5..262R。doi:10.3847/PSJ/ad86ba。
\item Marchi, S.;拉博尼，A.；Prettyman, T.H、; De Sanctis，M。C.;卡斯蒂略-罗杰兹，J.；雷蒙德，C。A.;Ammannito, E.;保龄球，T.；Ciarniello, M.;卡普兰，H.; 帕隆巴，E.；罗素，C。T.;维诺格拉多夫，V.；山下，N。(2018)。“一种水生改变的富含碳的谷神星”。自然天文学。3(2):140-145.doi:10.1038/s41550-018-0656-0。S2CID135013590。 
\item “在Ceres上更改了名称”。美国地质勘探局。2016年12月7日。已存档从原件到2021年8月19日。已检索8月19日2021。
\item 伊丽莎白·兰道 (2015年7月28日)。“Ceres的新名称和见解”。NASA。已存档从2016年1月6日起。已检索7月28日2015。
\item Marchi, S.;Ermakov, A.I.;雷蒙德，C。A.;傅，R。R、；O'Brien, D.P.;布兰德，M.T.;Ammannito, E.;De Sanctis, M.C.;保龄球，T.；申克，P.；Scully, J.E.C.;布茨科夫斯基，D。L.;威廉姆斯，D.A.;希辛格，H; 罗素，C。T.(2016年7月26日)。“谷神星上失踪的大型撞击坑”。自然通讯。712257。Bibcode:2016NatCo...712257M。doi:10.1038/ncomms12257。PMC4963536。PMID27459197。  
\item 威廉姆斯，大卫·A；Kneiss, T.(2018年12月)。“矮行星Ceres的Kerwan四边形的地质: 调查Ceres最古老，最大的撞击盆地”。伊卡洛斯。316:99-113。Bibcode:2018年Icar .. 316...99瓦。doi:10.1016/j.icarus.2017.08.015。S2CID85539501。已存档从原件到2021年8月16日。已检索8月16日2021。 
\item Nathues, A.;普拉茨，T.；唐卡，G.；霍夫曼，M.；史高丽，J.E.C.；斯坦因，N.；Ruesch, O.;门格尔，K。(2019年)。“最高空间分辨率下的彩色Occator陨石坑”。伊卡洛斯。320:24-38。Bibcode:2019年Icar .. 320...24N。doi:10.1016/j.icarus.2017.12.021。ISSN0019-1035。 
\item 斯特罗姆，R.G.；Marchi, S.;马尔霍特拉，R.(2018)。“谷神星和类地行星撞击陨石坑记录”。伊卡洛斯。302:104-108。arXiv:1804.01229。Bibcode:2018年Icar .. 302 .. 104s。doi:10.1016/j.icarus.2017.11.013。S2CID119009942。 
\item "谷神星上的Hanami Planum"。NASA.2018年3月23日。已存档从原件到2021年9月29日。已检索8月17日2021。
\item 施罗德，斯特凡·E; 卡森蒂，乌里; 豪伯，恩斯特; 雷蒙德，卡罗尔; 罗素，克里斯托弗 (2021年5月)。“矮行星谷神星的易碎巨石”。行星科学杂志。2(3): 111。arXiv:2105.11841。Bibcode:2021PSJ.....2..111S。doi:10.3847/PSJ/abfe66。S2CID235187212。 
\item 罗伯特·J·斯特恩；Gerya, Taras; Tackley, Paul J.(2018年1月)。“停滞的盖构造: 硅酸盐行星，矮行星，大卫星和大型小行星的观点”。地学前沿。9(1):103-119。Bibcode:2018GeoFr...9 .. 103S。doi:10.1016/j.gsf.2017.06.004。高密度脂蛋白:20.500.11850/224778。
\item “谷神星一次夺走了一座冰火山的生命”。亚利桑那大学。2018年9月17日。已存档从原件到2020年11月9日。已检索4月22日2019。
\item 布茨科夫斯基，D.；Scully, J.E.C.;雷蒙德，C。A.;罗素，C。T.(2017年12月)。“探索灶神星和谷神星的构造活动”。美国地球物理联盟，2017年秋季会议，摘要 # P53G-02。2017。Bibcode:2017AGUFM.P53G..02B。已存档从原件到2021年9月29日。已检索8月19日2021。
\item “PIA20348: 从拉莫看到的Ahuna Mons”。喷气推进实验室。2016年3月7日。已存档从2016年3月11日的原件。已检索4月14日2016。
\item 迈克尔·T·索里；Sizemore，Hanna G.等人 (2018年12月)。“地形揭示的谷神星上的低温火山速率”。自然天文学。2(12):946-950。Bibcode:2018NatAs...2 .. 946S。doi:10.1038/s41550-018-0574-1。S2CID186800298。已检索8月17日2021。 
\item Ruesch, O.;普拉茨，T.；申克，P.；麦克法登，L.A.;卡斯蒂略-罗盖兹，J.C.;快，L.C.;伯恩，S.；Preusker, F.;OBrien, D.P.;Schmedemann, N.;威廉姆斯，D.A.;李，J.- Y.；布兰德，M.T.;希辛格，H.; 克奈斯尔，T.；尼斯曼，A.；谢弗，M.；帕斯卡特，J.施密特，B。E.;布茨科夫斯基，D。L.;Sykes, M.五、；Nathues, A.;Roatsch, T.;霍夫曼，M.；雷蒙德，C。A.;罗素，C。T.(2016年9月2日)。“谷神星上的低温火山活动”。科学。353(6303) aaf4286。Bibcode:2016Sci...353.4286R。doi:10.1126/科学.aaf4286。PMID27701087。 
\item 迈克尔·M·索里；伯恩，谢恩; 平淡无奇，迈克尔·T；阿里·M·布拉姆森；安东一世·埃尔马科夫；克里斯托弗·W·汉密尔顿；凯瑟琳娜·奥托；鲁埃什，奥塔维亚诺; 罗素，克里斯托弗·T。(2017)。“消失的谷神星低温火山”(PDF)。地球物理研究快报。44(3):1243-1250。Bibcode:2017乔治。。44.1243s。doi:10.1002/2016GL072319。高密度脂蛋白:10150/623032。S2CID52832191。已存档从原件到2021年9月29日。已检索8月25日2019。  
\item “新闻-Ceres景点在最新的Dawn图像中继续神秘化”。NASA/JPL。已存档从原件到2021年7月25日。已检索7月25日2021。
\item “美国地质调查局: 谷神星命名法”(PDF)。已存档(PDF)从2015年11月15日的原件。已检索7月16日2015。
\item "Cerealia Facula"。行星命名法地名。USGS天体地质研究计划。
\item "Vinalia Faculae"。行星命名法地名。USGS天体地质研究计划。
\item 兰道，伊丽莎白; 麦卡特尼，格雷琴 (2018年7月24日)。“地球上什么看起来像谷神星？”。NASA。已存档从原件到2021年5月31日。已检索7月26日2021。
\item Schenk，Paul; Sizemore，Hanna; 等人 (2019年3月1日)。“谷神星Occator火山口的Cerealia Facula明亮沉积物和底部沉积物的中央坑和圆顶: 形态，比较和形成”。伊卡洛斯。320:159-187.Bibcode:2019年Icar .. 320 .. 159S。doi:10.1016/j.icarus.2018.08.010。S2CID125527752。 
\item 安德鲁·里夫金 (2015年7月21日)。“谷神星的黎明: Occator火山口的阴霾？”。行星社会。已存档从2016年5月14日开始。已检索3月8日2017。
\item 雷德，诺拉·泰勒。“谷神星上的水冰增加了埋藏海洋的希望 [视频]”。科学美国人。已存档从2016年4月7日的原件。已检索4月7日2016。
\item 伊丽莎白·兰道 (2015年12月9日)。“Ceres的亮点和起源的新线索”。phys.org。已存档从2015年12月9日的原件。已检索12月10日2015。
\item Vu, Tuan H.; Hodyss, Robert; Johnson, Paul V.;Choukroun，Mathieu (2017年7月)。“从冷冻的钠-氯化铵-碳酸盐盐水中优先形成钠盐-对ceres亮点的影响”。行星与空间科学。141:73-77.Bibcode:2017P & SS .. 141...73V。doi:10.1016/j.pss.2017.04.014。
\item 托马斯·B·麦考德；赞汶，弗朗西斯卡 (2019)。“黎明任务中谷神星的表面成分”。伊卡洛斯。318:2-13.Bibcode:2019年Icar .. 318。2米。doi:10.1016/j.icarus.2018.03.004。S2CID125115208。 
\item 快，Lynnae C.；黛布拉·L·布茨科夫斯基；鲁斯，奥塔维亚诺; 史高丽，詹妮弗·E。C.;朱莉·卡斯蒂略-罗杰兹; 卡罗尔·A·雷蒙德；保罗·M·申克；汉娜·G·西兹莫尔；Sykes, Mark V.(2019年3月1日)。“Occator火山口下方可能的盐水储层: 热和成分演化以及Cerealia Dome和Vinalia Faculae的形成”。伊卡洛斯。320:119-135.Bibcode:2019Icar .. 320 .. 119Q。doi:10.1016/j.icarus.2018.07.016。S2CID125508484。已存档从原件到2021年9月29日。已检索6月9日2021。 
\item 斯坦，N.T.;Ehlmann, B.L.(2019年3月1日)。“谷神星上亮点的形成和演变”。伊卡洛斯。320:188-201.Bibcode:2019年Icar .. 320 .. 188s。doi:10.1016/j.icarus.2017.10.014。
\item 麦卡特尼，格雷琴 (2020年8月11日)。“谜团解决了: 谷神星上的明亮区域来自下面的咸水”。Phys.org。已存档从原件到2020年8月11日。已检索8月12日2020。
\item 布兰德，迈克尔·T；雷蒙德，卡罗尔·A; 等人 (2016)。“陨石坑形态揭示的谷神星浅层地下的组成和结构”。自然地球科学。9(7):538-542。Bibcode:2016NatGe...9 .. 538B。doi:10.1038/ngeo2743。高密度脂蛋白:10919/103024。已存档从原件到2021年9月15日。已检索9月15日2021。
\item “PIA22660: ceres的内部结构 (艺术家的概念)”。照片日记帐。喷气推进实验室。2018年8月14日。已存档从原件到2019年4月21日。已检索4月22日2019。 公共领域本文结合了此来源的文本，该文本位于公共领域。
\item 奈夫，M.；Desch, S.J.(2016)。“带有泥泞冰幔的谷神星上的地球化学，热演化和低温火山作用”。地球物理研究快报。42(23)。doi:10.1002/2015GL066375。S2CID51756619。 
\item “确认: 谷神星有一个短暂的气氛”。今天的宇宙。2017年4月6日。已存档从原始的2017年4月15日。已检索4月14日2017。
\item Küppers, M.;奥罗克，L.；Bockelée-Morvan, D.；扎哈罗夫，V.；李，S.；冯·奥尔曼，P.；携带，B.；Teyssier, D.;马斯顿，A.；Müller, T.;克罗维西耶，J.；巴鲁奇，M.A.;莫雷诺，R。(2014年1月23日)“矮行星 (1) 谷神星上水蒸气的局部来源”。自然。505(7484):525-527。Bibcode:2014 Natur.505..525K。doi:10.1038/nature12918。ISSN0028-0836。PMID24451541。S2CID4448395。   
\item 坎平斯，H.; 舒适，C。M。(2014年1月23日)“太阳系: 蒸发小行星”。自然。505(7484):487-488。Bibcode:2014 Natur.505..487C。doi:10.1038/505487a。PMID24451536。S2CID4396841。  
\item 汉森，C。J.;埃斯波西托，L.；斯图尔特，A。I.;Colwell, J.;亨德里克斯，A.；普赖尔，W.；谢曼斯基，D.；西，R。(2006年3月10日)。“土卫二的水蒸气羽流”。科学。311(5766):1422-1425.Bibcode:2006Sci...311.1422H。doi:10.1126/science.1121254。PMID16527971。S2CID2954801。  
\item 罗斯，L.；萨尔，J.；雷瑟福德，K。D.;斯特罗贝尔，D。F.;费尔德曼，P。D.;麦格拉思，M。A.;尼莫，F。(2013年11月26日)。“欧罗巴南极的瞬态水蒸气”(PDF)。科学。343(6167):171-174.Bibcode:2014Sci...343 .. 171R。doi:10.1126/science.1247051。PMID24336567。S2CID27428538。已存档(PDF)从2013年12月16日的原件。已检索1月26日2014。   
\item O'Brien, D.P.;特拉维斯，B。J.;费尔德曼，W.C.;Sykes, M.五、；Schenk, P.M、；Marchi, S.;罗素，C。T.;雷蒙德，C。A.(2015年3月)。“由于地壳增厚和地下海洋加压，谷神星上有可能发生火山活动”(PDF)。46号月球和行星科学会议。第2831页。已存档(PDF)从2016年11月5日的原件。已检索3月1日2015。
\item Jewitt，David; Hsieh，Henry; Agarwal，Jessica (2015)。“活跃的小行星”(PDF)。见Michel，P.; 等人 (编辑)。小行星IV。亚利桑那大学。pp。 221-241.arXiv:1502.02361。Bibcode:2015aste.book .. 221J。doi:10.2458/azu_uapress_9780816532131-ch012。ISBN  978-0-8165-3213-1。S2CID 119209764。已存档 (PDF)从原件到2021年8月30日。已检索1月30日2020。
\item Jewitt, D; Chizmadia, L.;格林，R.；Prialnik，D (2007年)。“太阳系小天体中的水”(PDF)。在Reipurth，B.；Jewitt, D.;凯尔，K. (编辑)。原恒星和行星V。亚利桑那大学出版社。pp. 863-878。ISBN  978-0-8165-2654-3。已存档 (PDF)从2017年8月10日的原件。已检索10月11日2012。
\item 托马斯·B·麦考德；让-菲利普·科姆; 朱莉·C·卡斯蒂略·罗杰斯；哈里·Y·麦克斯温；Prettyman，托马斯·H (2022年5月)。谷神星，一个潮湿的星球: 黎明后的景色。地球化学。82(2) 125745。Bibcode:2022ChEG...82l5745M。doi:10.1016/j.chemer.2021.125745。
\item 希辛格，H.; 马尔基，S.；Schmedemann, N.;申克，P.；帕斯卡特，J.H、; Neesemann，A.；OBrien, D.P.;Kneissl, T.;Ermakov, A.I.;傅，R。R、；布兰德，M.T.;Nathues, A.;普拉茨，T.；威廉姆斯，D.A.;Jaumann, R.;卡斯蒂略-罗盖兹，J.C.;Ruesch, O.;施密特，B.；公园，R。美国；Preusker, F.;布茨科夫斯基，D。L.;罗素，C。T.;雷蒙德，C。A.(2016年9月1日)。“谷神星上的陨石坑: 对其地壳和进化的影响”。科学。353(6303) aaf4759。Bibcode:2016Sci...353.4759H。doi:10.1126/科学.aaf4759。PMID27701089。 
\item NASA/喷气推进实验室 (2016年9月1日)。“谷神星的地质活动，新研究揭示了冰”。科学日报。已存档从2017年4月5日开始。已检索3月8日2017。
\item 罗素，C。T.;雷蒙德，C。A.;Ammannito, E.;布茨科夫斯基，D。L.;De Sanctis, M.C.;希辛格，H.; Jaumann，R.；Konopliv, A.美国；麦克斯温，H. Y.；Nathues, A.;公园，R。S。(2016年9月2日)。“黎明到达谷神星: 探索一个小而多变的世界”。科学。353(6303):1008-1010.Bibcode:2016Sci...353.1008R。doi:10.1126/科学.aaf4219。ISSN0036-8075。PMID27701107。S2CID33455833。   
\item 诺伯特·舒尔霍弗; 谢恩·伯恩; 玛格丽特·E·兰迪斯；马扎里科，埃尔万; Prettyman，托马斯·H; 施密特，布兰妮·E；米凯拉·比利亚雷亚尔；朱莉·卡斯蒂略-罗杰兹; 卡罗尔·A·雷蒙德；罗素克里斯托弗·T.(2017年11月20日)。“假定的Cerean外逸层”。天体物理学杂志。850(1): 85。Bibcode:2017ApJ...850...85S。doi:10.3847/1538-4357/aa932f。高密度脂蛋白:10150/626261。ISSN0004-637X。 
\item 舍格霍夫，诺伯特; 本纳，迈赫迪; 贝雷日诺伊，阿列克谢A.；本杰明·格林哈根; 布兰特·M·琼斯；李帅; 奥兰多，托马斯·M；Prem, Parvathy; Tucker, Orenthal J.;W ö hler，基督徒 (2021年9月1日)。“月球，水星和谷神星上的水团外层和表面相互作用”。空间科学评论。217(6): 74。Bibcode:2021年SSRv ..。217。74s。doi:10.1007/s11214-021-00846-3。ISSN1572-9672。 
\item K ü ppers，迈克尔; O'Rourke，Laurence; Bockelée-morvan，Dominique; Zakharov，Vladimir; Lee，Seungwon; von Allmen，Paul; Carry，beno î t; Teyssier，David; Marston，Anthony；托马斯·穆勒; 雅克·克罗维西耶; 巴鲁奇，M。安东尼埃塔; 莫雷诺，拉斐尔 (2014年1月)。“矮行星 (1) 谷神星上水蒸气的局部来源”。自然。505(7484):525-527。Bibcode:2014 Natur.505..525K。doi:10.1038/nature12918。ISSN1476-4687。PMID24451541。  
\item Prettyman, T.H、; N.山下；Toplis, M.J.;麦克斯温，H. Y.；Schörghofer, N.;Marchi, S.;费尔德曼，W.C.;卡斯蒂略-罗杰兹，J.；Forni, O.;劳伦斯，D。J.;Ammannito, E.;Ehlmann, B.L.;Sizemore, H. G.;Joy, S.P.;波兰斯基，C。A.(2017年1月6日)。“谷神星的水状蚀变岩中的大量水冰: 来自核光谱学的证据”。科学。355(6320):55-59。Bibcode:2017Sci...355...55P。doi:10.1126/科学.aah6765。ISSN0036-8075。PMID27980087。  
\item 安德鲁·S·里夫金；李健阳; 米利肯，拉尔夫·E；Lim, Lucy F.;艾米·J·洛弗尔；布兰妮·E·施密特；露西·麦克法登；芭芭拉·科恩.(2011年12月1日)。“谷神星的表面组成”。空间科学评论。163(1):95-116。Bibcode:2011SSRv。。163...95R。doi:10.1007/s11214-010-9677-4。ISSN1572-9672。 
\item 涂，L.；Ip, W.-H.;王，Y。-C.(2014年12月1日)。“一个升华驱动的谷神星外层模型”。行星与空间科学。104:157-162。Bibcode:2014年P & SS .. 104 .. 157吨。doi:10.1016/j.pss.2014.09.002。ISSN0032-0633。 
\item 海恩，P。O.;Aharonson，O。(2015年9月)。“具有粗糙地形的谷神星上的冰的热稳定性”。地球物理研究杂志: 行星。120(9):1567-1584。Bibcode:2015JGRE..120.1567H。doi:10.1002/2015JE004887。ISSN2169-9097。 
\item 托马斯·B·麦考德；露西·麦克法登；克里斯托弗·T·罗素；索丁，克里斯托夫; 托马斯，彼得C。(2006年3月7日)。“谷神星，灶神星和帕拉斯: 原行星，而不是小行星”。Eos。87(10): 105。Bibcode:2006年EOSTr。。87 .. 105米。doi:10.1029/2006EO100002。已存档从原件到2021年9月28日。已检索9月12日2021。
\item 杨继金; 戈德斯坦，约瑟夫一世; 斯科特，爱德华·R。D.(2007)。“铁陨石早期形成和原行星灾难性破坏的证据”。自然。446(7138):888-891。Bibcode:2007 Natur.446..888Y。doi:10.1038/nature05735。PMID17443181。S2CID4335070。已存档从原件到2021年9月29日。已检索9月16日2021。  
\item 佩蒂特，让-马克; 莫比德利，亚历山德罗 (2001年)。“小行星带的原始激发和清除”(PDF)。伊卡洛斯。153(2):338-347.Bibcode:2001年Icar .. 153 .. 338页。doi:10.1006/icar.2001.6702。已存档(PDF)从2007年2月21日的原件。已检索6月25日2009年。
\item Greicius，Tony (2016年6月29日)。“最近的热液活动可能解释了ceres最明亮的区域”。nasa.gov。已存档从原始的2019年1月6日。已检索7月26日2016。
\item 南希阿特金森 (2016年7月26日)。“谷神星上的大型撞击坑不见了”。今天的宇宙已存档从原件到2021年5月15日。已检索5月15日2021。
\item 墙，迈克 (2016年9月2日)。“NASA的黎明任务在谷神星上监视冰火山”。科学美国人。已存档从原件到2017年6月3日。已检索3月8日2017。
\item 卡斯蒂略-罗盖兹，J.C.;麦考德，T.B.;戴维斯，A。G.(2007)。“谷神星: 进化和现状”(PDF)。月球和行星科学。三十八:2006-2007。已存档(PDF)从2011年2月24日的原件。已检索6月25日2009年。
\item De Sanctis, M.C.;维诺格拉多夫，V.；拉博尼，A.；Ammannito, E.;Ciarniello, M.;Carrozzo, F.G.;De Angelis, S.;雷蒙德，C。A.;罗素，C。T.(2018年10月17日)。“来自VIR/Dawn高空间分辨率光谱的Ceres上的有机物特征”。皇家天文学会月刊。482(2):2407-2421。doi:10.1093/mnras/sty2772。
\item 布兰登·斯佩克特 (2021年1月19日)。“天体物理学家说，人类可能会在未来15年内搬到这个漂浮的小行星带殖民地。”。活的科学。已存档从原件到2021年6月24日。已检索6月23日2021。
\item 门泽尔，唐纳德·H; 帕萨乔夫，杰伊·M。(1983)。恒星和行星的野外指南(第二版)。波士顿:霍顿·米夫林。p。 391。ISBN 978-0-395-34835-2。
\item 马丁内斯，帕特里克 (1994)。观察者天文学指南。剑桥大学出版社。第298页。ISBN 978-0-521-37945-8。OCLC 984418486。
\item 米利斯，L.R、；Wasserman, L.H、; 弗朗兹，O。Z.; 等人 (1987年)。“谷神星的大小，形状，密度和反照率来自其对BD 8 ° 471的掩星”。伊卡洛斯。72(3):507-518。Bibcode:1987年伊卡尔。72 .. 507米。doi:10.1016/0019-1035(87)90048-0。高密度脂蛋白:2060/19860021993。
\item “凯克自适应光学成像矮行星谷神星”。自适应光学。2006年10月11日。存档自原来的2009年8月18日。已检索4月27日2007。
\item 最大的小行星可能是带有水冰的 “迷你行星”。HubbleSite。2005年9月7日。已存档从2021年7月20日开始。已检索7月20日2021。
\item Carry，Benoit; 等人 (2007)。“矮行星谷神星的近红外测绘和物理特性”(PDF)。天文学与天体物理学。478(1):235-244。arXiv:0711.1152。Bibcode:2008A & A...478 .. 235C。doi:10.1051/0004-6361:20078166。S2CID6723533。存档自原来的(PDF)2008年5月30日。  
\item Houtkooper, J.M、；Schulze-Makuch, D.(2017)。“谷神星: 天体生物学的前沿”(PDF)。天体生物学科学会议(1965)。已存档(PDF)从原件到2021年8月30日。已检索8月19日2021。
\item 罗素，C。T.;卡帕乔尼，F.；Coradini，A.; 等人 (2007年10月)。“黎明任务到Vesta和谷神星”(PDF)。地球、月球和行星。101(1-2):65-91。Bibcode:2007EM & P。。101...65R。doi:10.1007/s11038-007-9151-9。S2CID46423305。已存档(PDF)从原件到2020年10月25日。已检索6月13日2011年。  
\item 库克，佳瑞C.；布朗，德韦恩C.(2011年5月11日)。“NASA的黎明捕捉到第一张接近小行星的图像”。NASA/JPL。已存档从2011年5月14日的原件。已检索5月14日2011年。
\item Schenk, P.(2015年1月15日)。“矮人年”: 谷神星和冥王星应得的。行星社会。已存档从原件到2015年2月21日。已检索2月10日2015。
\item Rayman，Marc (2014年12月1日)。“黎明日报: 展望谷神星”。行星社会。已存档从原件到2015年2月26日。已检索3月2日2015。
\item 罗素，C。T.;卡帕乔尼，F.；Coradini，A.; 等人 (2006年)。“黎明发现任务到灶神星和谷神星: 现状”。空间研究进展。38(9):2043-2048。arXiv:1509.05683。Bibcode:2006年AdSpR。。38.2043r。doi:10.1016/j.asr.2004.12.041。
\item Rayman，Marc (2015年1月30日)。“黎明日报: 关闭谷神星”。行星社会。已存档从2015年3月1日起。已检索3月2日2015。
\item 雷曼，马克 (2015年3月6日)。“黎明日记: 谷神星轨道插入!”。行星社会。已存档从原来的2015年3月8日。已检索3月6日2015。
\item Rayman，Marc (2014年3月3日)。“黎明日记: 在谷神星周围机动”。行星社会。已存档从原件到2015年2月26日。已检索3月6日2015。
\item Rayman，Marc (2014年4月30日)。“黎明杂志: 解释轨道插入”。行星社会。已存档从原件到2015年2月26日。已检索3月6日2015。
\item 雷曼，马克 (2014年6月30日)。“黎明日报: 谷神星的哈莫”。行星社会。已存档从原件到2015年2月26日。已检索3月6日2015。
\item Rayman，Marc (2014年8月31日)。“黎明日报: 从哈莫到拉莫及以后”。行星社会。已存档从2015年3月1日起。已检索3月6日2015。
\item “来自Ceres的黎明数据公开发布: 最后，为全球肖像着色!”。行星社会。已存档从2015年11月9日的原件。已检索11月9日2015。
\item “黎明任务在谷神星扩展”。NASA/JPL-加州理工学院。2017年10月19日。已存档从原件到2021年10月1日。已检索10月1日2021。
\item 菲尔，辫子(2015年5月11日)。“谷神星的亮点旋转进入视野”。板岩。已存档从原件到2015年5月29日。已检索5月30日2015。
\item 奥尼尔，伊恩 (2015年2月25日)。“谷神星的神秘亮点可能有火山起源”。发现公司.已存档从2016年8月14日的原件。已检索3月1日2015。
\item 艾米丽·拉克达瓦拉(2015)。“LPSC 2015: 谷神星黎明的第一个结果: 临时地名和可能的羽毛”。行星社会。已存档从2016年5月6日的原件。已检索9月23日2021。
\item "Ceres RC3动画"。喷气推进实验室。2015年5月11日已存档从原件到2021年1月17日。已检索7月31日2015。
\item De Sanctis, M.C.; 等人 (2016年6月29日)。“明亮的碳酸盐沉积物是 (1) 谷神星含水蚀变的证据”。自然。536(7614):54-57.Bibcode:2016自然536. .. 54D。doi:10.1038/自然18290。PMID27362221。S2CID4465999。  
\item Rayman，Marc (2018年6月13日)。“黎明-任务状态”。喷气推进实验室。已存档从原件到2018年6月23日。已检索6月16日2018。
\item  雷曼，马克 (2018)。"亲爱的Dawntasmagorias"。NASA喷气推进实验室。已存档从原件到2021年7月21日。已检索7月21日2021。
\item 基西克，L.E.;Acciarini, G.;贝茨，H.; 等人 (2020年)。“从遗迹海洋世界返回的样本: 到谷神星Occator陨石坑的卡尔萨斯任务”(PDF)。第51届月球与行星科学会议。已存档(PDF)从原件到2020年10月26日。已检索2月1日2020。
\item 邹永辽; 李; 欧阳自远。“中国的深空探测到2030年”(PDF)。中国科学院月球与深空探测重点实验室，国家天文台，北京.已存档(PDF)从2014年12月14日的原件。已检索9月23日2021。
\end{enumerate}
\subsection{外部链接}
\begin{itemize}
\item Ceres Trek-1 Ceres的数据集和地图的集成地图浏览器
\item Ceres 3D模型-NASA
\item 目的地谷神星: 黎明早餐-NASA
\item 黎明任务主页在JPL
\item 谷神星轨道的模拟
\item Google Ceres 3D,矮行星的交互式地图
\item 高斯如何确定谷神星的轨道 已存档2008年4月14日在回程机fro m keplersdiscovery.com
\item 谷神星的动画重投影彩色地图(2015年2月22日)
\item 视频 (3:34): 谷神星 “亮点” -- 谜团解开 (2020年8月10日)开YouTube
\item 旋转浮雕模型塞恩·多兰的谷神星 (大约是全程旋转的60\%; 从中心上方的中途岛开始)
\item 谷神星 (矮行星)在AstDyS-2，小行星-动态站点
\item 星历 · 观测预测 · 轨道信息 · 适当的元素 · 观测信息
\item 谷神星在JPL小型数据库 在Wikidata上编辑此内容
\item 关闭进近 · 发现 · 星历 · 轨道查看器 · 轨道参数 · 物理参数
\end{itemize}
